%% kapa: exam=1 round=32 year=2021 period=1 form=A subject=civil_law subject_ko=민법 seq=2 number=2
%% source: Q-Net 감정평가사 기출문제 (cst00309) artlSeq=5208397 file=extracted/1008/5208397/A형/2021년도 제32회 감정평가사 1차 1교시 A형.pdf
%% conversion: kordoc 4.24.51 -> tools/gen1.py (auto); review=proofread-sonnet; flags=none
\begin{question}{2}
신의성실의 원칙에 관한 설명으로 옳은 것을 모두 고른 것은? (다툼이 있으면 판례에 따름)
\begin{framedtext}
ㄱ. 회사의 이사가 회사의 확정채무를 보증한 경우에는 그 직을 사임하더라도 사정변경을 이유로 그 보증계약을 해지할 수 없다. \par
ㄴ. 소멸시효 완성 전에 채무자가 시효중단을 현저히 곤란하게 하여 채권자가 아무런 조치를 취할 수 없었던 경우, 그 채무자가 시효완성을 주장하는 것은 신의칙 상 허용되지 않는다. \par
ㄷ. 강행법규를 위반한 자가 스스로 강행법규 위반을 이유로 약정의 무효를 주장하는 것은 특별한 사정이 없는 한 신의칙에 반한다. \par
\end{framedtext}
\begin{choices}
  \item ㄱ
  \item ㄷ
  \item ㄱ, ㄴ
  \item ㄴ, ㄷ
  \item ㄱ, ㄴ, ㄷ
\end{choices}
\end{question}
\answer{3}
